\chapter*{Abstract}
\addcontentsline{toc}{chapter}{Abstract}
\chaptermark{Abstract}

This thesis describes the design and development of the \textbf{SMART-SOJA} project, a mobile, autonomous, and connected soybean pre-treatment unit, developed to meet the strict quality standards of the Glo-Djigbé Industrial Zone (GDIZ) in Benin. Carried out as part of a graduation internship at \textbf{QOTTO Benin} for a Professional License in Industrial IT and Maintenance (EGEI/UCAO), this project addresses the limitations of traditional post-harvest methods that fail to maintain purity levels above 98\% and moisture levels below 12\% required by industrial crushing plants.

The methodology combines mechanical CAD modeling (SolidWorks), electronic circuit design (KiCad), real-time control (FreeRTOS on an ESP32 microcontroller), and Internet of Things (IoT) connectivity. The modular prototype incorporates a vibrating sieving system driven by a DC motor, an active forced-air drying chamber monitored by capacitive humidity sensors, a digital weighing module (load cell and HX711 amplifier), and a standalone 12 V battery powered by a solar photovoltaic panel.

Laboratory tests have validated the correct functioning of all subsystems under unit testing conditions: DHT22 temperature/humidity sensor, soybean moisture sensors, drive motors, heating element, ventilation fan, and LCD interface. MQTT data transmission achieves 100\% transmission integrity with no message loss. The system automatically generates a unique QR Code (``digital passport'') for each batch, ensuring robust traceability throughout the processing chain. However, the prototype represents a functional proof of concept at this stage. Validation under real-world loads (2 to 5 kg), precise evaluation of the complete cleaning + drying cycle (achievement of target purity and moisture rates), and final integration of the weighing module remain to be accomplished before operational deployment.

\textbf{Keywords:} Soybean, Pre-treatment, Industrial IT, Internet of Things, Traceability, Solar energy, GDIZ.
